\section{Linear Maps}

\subsection{Vector Space of Linear Maps}

\begin{exercise}[3a1]
    Suppose $b, c \in \r$. Define $T: \r^3 \to \r^2$ by
    \[T(x, y, z) = (2x - 4y + 3z + b, 6x + cxyz).\]
    Show that $T$ is linear if and only if $b = c = 0$.
\end{exercise}

\begin{sol}
    Suppose $T$ is linear. Then
    \[T(1, 1, 1) = T(1, 0, 0) + T(0, 1, 0) + T(0, 0, 1) = (2 + b, 6) + (-4 + b, 0) + (3 + b, 0) = (1 + 3b, 6).\]
    On the other hand,
    \[T(1, 1, 1) = (2 - 4 + 3 + b, 6 + c) = (1 + b, 6 + c).\]
    Comparing the two expressions, we get $1 + 3b = 1 + b$ and $6 = 6 + c$, which implies $b = 0$ and $c = 0$.

    Conversely, if $b = c = 0$, then $T(x, y, z) = (2x - 4y + 3z, 6x)$. Let $(x_1, y_1, z_1), (x_2, y_2, z_2) \in \r^3$ and $\alpha \in \r$. Then
    \begin{align*}
        T(x_1 + x_2, y_1 + y_2, z_1 + z_2) & = (2(x_1 + x_2) - 4(y_1 + y_2) + 3(z_1 + z_2), 6(x_1 + x_2)) \\
        & = (2x_1 + 2x_2 - 4y_1 - 4y_2 + 3z_1 + 3z_2, 6x_1 + 6x_2) \\
        & = (2x_1 - 4y_1 + 3z_1, 6x_1) + (2x_2 - 4y_2 + 3z_2, 6x_2) \\
        & = T(x_1, y_1, z_1) + T(x_2, y_2, z_2),
    \end{align*}
    and
    \begin{align*}
        T(\alpha x_1, \alpha y_1, \alpha z_1) & = (2(\alpha x_1) - 4(\alpha y_1) + 3(\alpha z_1), 6(\alpha x_1)) \\
        & = (\alpha(2x_1 - 4y_1 + 3z_1), \alpha(6x_1)) \\
        & = \alpha(2x_1 - 4y_1 + 3z_1, 6x_1) \\
        & = \alpha T(x_1, y_1, z_1).
    \end{align*}
    Therefore, $T$ is linear if and only if $b = c = 0$.
\end{sol}

\begin{exercise}[3a2]
    Suppose $b, c \in \r$. Define $T: \pp(\r) \to \r^2$ by
    \[Tp = \left(3p(4) + 5p'(6) + bp(1)p(2), \int_{-1}^{2} x^3 p(x) \dd x + c \sin p(0)\right).\]
    Show that $T$ is linear if and only if $b = c = 0$.
\end{exercise}

\begin{sol}
    Suppose $T$ is linear. Then for any $p \in \pp(\r)$, we have $T(2p) = 2Tp$. Looking at the first component of $T(2p)$, we have
    \[3(2p)(4) + 5(2p)'(6) + b(2p)(1)(2p)(2) = 6p(4) + 10p'(6) + 4bp(1)p(2).\]
    On the other hand, the first component of $2Tp$ is
    \[2(3p(4) + 5p'(6) + bp(1)p(2)) = 6p(4) + 10p'(6) + 2bp(1)p(2).\]
    Comparing the two expressions, we get $4bp(1)p(2) = 2bp(1)p(2)$ for all $p \in \pp(\r)$. Since there exist polynomials $p$ such that $p(1)p(2) \neq 0$, we must have $b = 0$.

    The second component of $T(2p)$ must be equal to the second component of $2Tp$. We have
    \[\int_{-1}^{2} x^3 (2p)(x) \dd x + c \sin(2p(0)) = 2\int_{-1}^{2} x^3 p(x) \dd x + 2c \sin (p(0)).\]
    The integral terms may be subtracted, so we get $c \sin(2p(0)) = 2c \sin(p(0))$ for all $p \in \pp(\r)$. If $c \neq 0$, then we can divide both sides by $c$ to get $\sin(2p(0)) = 2\sin(p(0))$, which implies $\sin(2x) = 2\sin(x)$ for all $x \in \r$. Then $\cos(x) = 1$ for all $x \in \r$, which is a contradiction. Therefore, $c = 0$.

    Conversely, if $b = c = 0$, let $p, q \in \pp(\r)$ and $\alpha \in \r$. Then
    \begin{align*}
        T(p + q) & = \left(3(p + q)(4) + 5(p + q)'(6), \int_{-1}^{2} x^3 (p + q)(x) \dd x\right) \\
        & = \left(3p(4) + 3q(4) + 5p'(6) + 5q'(6), \int_{-1}^{2} x^3 p(x) \dd x + \int_{-1}^{2} x^3 q(x) \dd x\right) \\
        & = \left(3p(4) + 5p'(6), \int_{-1}^{2} x^3 p(x) \dd x\right) + \left(3q(4) + 5q'(6), \int_{-1}^{2} x^3 q(x) \dd x\right) \\
        & = Tp + Tq,
    \end{align*}
    and
    \begin{align*}
        T(\alpha p) & = \left(3(\alpha p)(4) + 5(\alpha p)'(6), \int_{-1}^{2} x^3 (\alpha p)(x) \dd x\right) \\
        & = \left(\alpha (3p(4) + 5p'(6)), \alpha \int_{-1}^{2} x^3 p(x) \dd x\right) \\
        & = \alpha\left(3p(4) + 5p'(6), \int_{-1}^{2} x^3 p(x) \dd x\right) \\
        & = \alpha Tp.
    \end{align*}
    Therefore, $T$ is linear if and only if $b = c = 0$.
\end{sol}

\begin{exercise}[3a3]
    Suppose that $T \in \lin{\f^n, \f^m}$. Show that there exist scalars $A_{j, k} \in \f$ for $j = 1, \ldots, m$ and $k = 1, \ldots, n$ such that
    \[T(x_1, \ldots, x_n) = (A_{1,1} x_1 + \cdots + A_{1, n} x_n, \ldots, A_{m, 1} x_1 + \cdots + A_{m, n} x_n)\]
    for every $(x_1, \ldots, x_n) \in \f^n$.
\end{exercise}

\begin{sol}
    Let $e_1, \ldots, e_n$ be the standard basis of $\f^n$. Then for each $j = 1, \ldots, m$, we have
    \[T(e_k) = (A_{1, k}, A_{2, k}, \ldots, A_{m, k})\]
    for some scalars $A_{j, k} \in \f$. Now let $(x_1, \ldots, x_n) \in \f^n$. Then
    \begin{align*}
        T(x_1, \ldots, x_n) & = T(x_1 e_1 + \cdots + x_n e_n) \\
        & = x_1 T(e_1) + \cdots + x_n T(e_n) \\
        & = x_1 (A_{1, 1}, A_{2, 1}, \ldots, A_{m, 1}) + \cdots + x_n (A_{1, n}, A_{2, n}, \ldots, A_{m, n}) \\
        & = (A_{1, 1} x_1 + \cdots + A_{1, n} x_n, A_{2, 1} x_1 + \cdots + A_{2, n} x_n, \ldots, A_{m, 1} x_1 + \cdots + A_{m, n} x_n).
    \end{align*}
    Therefore,
    \[T(x_1, \ldots, x_n) = (A_{1, 1} x_1 + \cdots + A_{1, n} x_n, A_{2, 1} x_1 + \cdots + A_{2, n} x_n, \ldots, A_{m, 1} x_1 + \cdots + A_{m, n} x_n). \qh\]
\end{sol}

\begin{exercise}[3a4]
    Suppose $T \in \lin{V, W}$ and $v_1, \ldots, v_m$ is a list of vectors in $V$ such that $Tv_1, \ldots, Tv_m$ is a linearly independent list in $W$. Prove that $v_1, \ldots, v_m$ is linearly independent.
\end{exercise}

\begin{sol}
    Suppose that $a_1 v_1 + \cdots + a_m v_m = 0$ for some scalars $a_1, \ldots, a_m \in \f$. Applying $T$ to both sides, we get
    \[a_1 Tv_1 + \cdots + a_m Tv_m = T(0) = 0.\]
    Since $Tv_1, \ldots, Tv_m$ is linearly independent, it follows that $a_1 = \cdots = a_m = 0$. Therefore, $v_1, \ldots, v_m$ is linearly independent.
\end{sol}

\begin{exercise}[3a5]
    Prove that $\lin{V, W}$ is a vector space, as was asserted in 3.6.
\end{exercise}

\begin{sol}
    We must first ensure that the operations of addition and scalar multiplication are well-defined, i.e., if $S, T \in \lin{V, W}$ and $\alpha \in \f$, then $S + T \in \lin{V, W}$ and $\alpha T \in \lin{V, W}$. To that end, let $v, u \in V$. Then
    \begin{align*}
        (S + T)(v + u) & = S(v + u) + T(v + u) \\
        & = (Sv + Su) + (Tv + Tu) \\
        & = (Sv + Tv) + (Su + Tu) \\
        & = (S + T)v + (S + T)u,
    \end{align*}
    and
    \begin{align*}
        (\alpha T)(v + u) & = \alpha T(v + u) \\
        & = \alpha (Tv + Tu) \\
        & = \alpha Tv + \alpha Tu \\
        & = (\alpha T)v + (\alpha T)u.
    \end{align*}
    Also, for any $v \in V$ and $\alpha \in \f$,
    \begin{align*}
        (S + T)(\alpha v) & = S(\alpha v) + T(\alpha v) \\
        & = \alpha Sv + \alpha Tv \\
        & = \alpha (Sv + Tv) \\
        & = \alpha (S + T)v,
    \end{align*}
    and
    \begin{align*}
        (\alpha T)(\beta v) & = \alpha T(\beta v) \\
        & = \alpha (\beta Tv) \\
        & = (\alpha \beta) Tv \\
        & = (\alpha \beta) T v.
    \end{align*}
    Therefore, $S + T \in \lin{V, W}$ and $\alpha T \in \lin{V, W}$. We now check the vector space axioms:
    \begin{enumerate}
        \item \textbf{Commutativity}: Let $S, T \in \lin{V, W}$. Then for any $v \in V$,
        \begin{align*}
            (S + T)v & = Sv + Tv \\
            & = Tv + Sv \\
            & = (T + S)v.
        \end{align*}
        \item \textbf{Associativity}: Let $R, S, T \in \lin{V, W}$. Then for any $v \in V$,
        \begin{align*}
            ((R + S) + T)v & = (R + S)v + Tv \\
            & = (Rv + Sv) + Tv \\
            & = Rv + (Sv + Tv) \\
            & = Rv + (S + T)v \\
            & = (R + (S + T))v.
        \end{align*}
        Similarly, for any $\alpha, \beta \in \f$ and $T \in \lin{V, W}$,
        \begin{align*}
            ((\alpha\beta)T)v & = (\alpha\beta)(Tv) \\
            & = \alpha(\beta Tv) \\
            & = \alpha(\beta T)v \\
            & = (\alpha(\beta T))v.
        \end{align*}
        \item \textbf{Additive identity}: Let $0 \in \lin{V, W}$ be the zero map defined by $0v = 0$ for all $v \in V$. Then for any $T \in \lin{V, W}$ and $v \in V$,
        \[(T + 0)v = Tv + 0v = Tv + 0 = Tv.\]
        \item \textbf{Inverse}: Let $T \in \lin{V, W}$. Define $-T \in \lin{V, W}$ by $(-T)v = -Tv$ for all $v \in V$. Then for any $v \in V$,
        \[(T + (-T))v = Tv + (-T)v = Tv - Tv = 0 = 0v.\]
        \item \textbf{Multiplicative identity}: Let $1 \in \f$ be the multiplicative identity. Then for any $T \in \lin{V, W}$ and $v \in V$,
        \[(1T)v = 1(Tv) = Tv.\]
        \item \textbf{Distributive properties}: Let $T, S \in \lin{V, W}$ and $\alpha, \beta \in \f$. Then for any $v \in V$,
        \begin{align*}
            (\alpha(T + S))v & = (\alpha T + \alpha S)v \\
            & = (\alpha T)v + (\alpha S)v \\
            & = \alpha(Tv) + \alpha(Sv) \\
            & = \alpha(Tv + Sv) \\
            & = \alpha((T + S)v),
        \end{align*}
        and
        \begin{align*}
            ((\alpha + \beta)T)v & = (\alpha + \beta)(Tv) \\
            & = \alpha(Tv) + \beta(Tv) \\
            & = (\alpha T)v + (\beta T)v \\
            & = (\alpha T + \beta T)v. \qh
        \end{align*}
    \end{enumerate}
\end{sol}

\begin{exercise}[3a6]
    Prove that multiplication of linear maps has the associative, identity, and distributive properties asserted in 3.8.
    \begin{subexercises}
        \item \textbf{Associativity}: $(T_1T_2)T_3 = T_1(T_2T_3)$ whenever $T_1, T_2,$ and $T_3$ are linear maps such that the products make sense (meaning $T_3$ maps into the domain of $T_2$, and $T_2$ maps into the domain of $T_1$).
        \item \textbf{Identity}: $TI = IT = T$ whenever $T \in \lin{V, W}$; here the first $I$ is the identity operator on $V$ and the second $I$ is the identity operator on $W$.
        \item \textbf{Distributive properties}: $(S_1 + S_2)T = S_1T + S_2T$ and $S(T_1 + T_2) = ST_1 + ST_2$ whenever $T, T_1, T_2 \in \lin{U, V}$ and $S, S_1, S_2 \in \lin{V, W}$.
    \end{subexercises}
\end{exercise}

\begin{sole}
    \item Let $T_1 \in \lin{V, W}$, $T_2 \in \lin{U, V}$, and $T_3 \in \lin{X, U}$. Then for any $x \in X$,
    \begin{align*}
        ((T_1T_2)T_3)x & = (T_1T_2)(T_3x) \\
        & = T_1(T_2(T_3x)) \\
        & = T_1((T_2T_3)x) \\
        & = (T_1(T_2T_3))x.
    \end{align*}
    \item Let $T \in \lin{V, W}$. Then for any $v \in V$,
    \[(TI)v = T(Iv) = Tv \quad \text{and} \quad (IT)v = I(Tv) = Tv.\]
    \item Let $T, T_1, T_2 \in \lin{U, V}$ and $S, S_1, S_2 \in \lin{V, W}$. Then for any $u \in U$,
    \begin{align*}
        ((S_1 + S_2)T)u & = (S_1 + S_2)(Tu) \\
        & = S_1(Tu) + S_2(Tu) \\
        & = (S_1T)u + (S_2T)u \\
        & = (S_1T + S_2T)u,
    \end{align*}
    and
    \begin{align*}
        (S(T_1 + T_2))u & = S((T_1 + T_2)u) \\
        & = S(T_1u + T_2u) \\
        & = S(T_1u) + S(T_2u) \\
        & = (ST_1)u + (ST_2)u \\
        & = (ST_1 + ST_2)u. \qh
    \end{align*}
\end{sole}

\begin{exercise}[3a7]
    Show that every linear map from a one-dimensional vector space to itself is multiplication by some scalar. More precisely, prove that if $\dim V = 1$ and $T \in \lin{V}$, then there exists $\lambda \in \f$ such that $T\varphi = \lambda\varphi$ for all $v \in V$.
\end{exercise}

\begin{sol}
    Suppose $\dim V = 1$ and let $v \in V$ be a nonzero vector. Then $\set{v}$ is a basis of $V$. Let $T \in \lin{V}$. Since $Tv \in V$, there exists $\lambda \in \f$ such that $Tv = \lambda v$. Now let $u \in V$. Then there exists $\alpha \in \f$ such that $u = \alpha v$. Therefore,
    \[Tu = T(\alpha v) = \alpha Tv = \alpha (\lambda v) = \lambda (\alpha v) = \lambda u. \qh\]
\end{sol}

\begin{exercise}[3a8]
    Give an example of a function $\varphi: \r^2 \to \r$ such that
    \[\varphi(av) = a\varphi(v)\]
    for all $a \in \r$ and all $v \in \r^2$ but $\varphi$ is not linear.

    \remark{This was a surprisingly difficult exercise for me, namely because there were no obvious examples of homogeneous functions that were not linear. However, the idea is that while $x^n$ functions aren't linear, they don't maintain homogeneity due to an extra $a^{n - 1}$ factor. It follows that we must divide $x^n$ by some $n - 1$ power to maintain homogeneity. However, dividing by just $y$ risks division by zero, so we must divide by $x + y$ to avoid this.}
\end{exercise}

\begin{sol}
    Let $\varphi: \r^2 \to \r$ be defined by
    \[\varphi(x, y) = \frac{x^2}{x + y}.\]
    Then for any $a \in \r$ and $(x, y) \in \r^2$,
    \[\varphi(a(x, y)) = \varphi(ax, ay) = \frac{(ax)^2}{ax + ay} = \frac{a^2 x^2}{a(x + y)} = a\frac{x^2}{x + y} = a\varphi(x, y).\]
    However, $\varphi$ is not additive. For example, let $v = (1, 0)$ and $w = (0, 1)$. Then
    \[\varphi(v + w) = \varphi(1, 1) = \frac{1^2}{1 + 1} = \frac{1}{2},\]
    but
    \[\varphi(v) + \varphi(w) = \varphi(1, 0) + \varphi(0, 1) = \frac{1^2}{1 + 0} + \frac{0^2}{0 + 1} = 1 + 0 = 1.\]
    Therefore, $\varphi$ is not linear.
\end{sol}

\begin{exercise}[3a9]
    Give an example of a function $\varphi: \c \to \c$ such that
    \[\varphi(w + z) = \varphi(w) + \varphi(z)\]
    for all $w, z \in \c$ but $\varphi$ is not linear. (Here $\c$ is thought of as a complex vector space.)
\end{exercise}

\begin{sol}
    Let $\varphi: \c \to \c$ be defined by
    \[\varphi(x + iy) = x.\]
    Then for any $w = x_1 + iy_1$ and $z = x_2 + iy_2$ in $\c$,
    \[\varphi(w + z) = \varphi((x_1 + x_2) + i(y_1 + y_2)) = x_1 + x_2 = \varphi(w) + \varphi(z).\]
    However, $\varphi$ is not homogeneous. For example, let $a = i$ and $z = 1$. Then
    \[\varphi(iz) = \varphi(i) = 0,\]
    but
    \[i\varphi(z) = i\varphi(1) = i(1) = i.\]
    Therefore, $\varphi$ is not linear.
\end{sol}

\begin{exercise}[3a10]
    Prove or give a counterexample: If $q \in \pp(\r)$ and $T: \pp(\r) \to \pp(\r)$ is defined by $Tp = q \circ p$, then $T$ is a linear map.
\end{exercise}

\begin{sol}
    This is false. Let $q(x) = x^2$ and define $T: \pp(\r) \to \pp(\r)$ by $Tp = q \circ p$. Then for any $p_1, p_2 \in \pp(\r)$,
    \[T(p_1 + p_2) = q \circ (p_1 + p_2) = (p_1 + p_2)^2 = p_1^2 + 2p_1p_2 + p_2^2 \neq p_1^2 + p_2^2 = Tp_1 + Tp_2. \qh\]
\end{sol}

\begin{exercise}[3a11]
    Suppose $V$ is finite-dimensional and $T \in \lin{V}$. Prove that $T$ is a scalar multiple of the identity if and only if $ST = TS$ for every $S \in \lin{V}$.

    \remark{The backwards direction is difficult. The idea is to construct many such $S$ that will force $T$ to be a scalar multiple of the identity.}
\end{exercise}

\begin{sol}
    Suppose $T$ is a scalar multiple of the identity, i.e., $T = \lambda I$ for some $\lambda \in \f$. Then for any $S \in \lin{V}$,
    \[ST = S(\lambda I) = \lambda S = (\lambda I)S = TS.\]

    Conversely, suppose $ST = TS$ for every $S \in \lin{V}$. Let $v_1, \ldots, v_n$ be a basis of $V$. Define $S_{ij} \in \lin{V}$ by
    \[S_{ij}v_k = 
    \begin{cases}
        v_j & \text{if } k = i, \\
        0 & \text{if } k \neq i.
    \end{cases}.\]
    In other words, $S_{ij}$ sends $v_i$ to $v_j$ and sends all other basis vectors to $0$. Fix some $i$, and let $Tv_i = \alpha_1 v_1 + \cdots + \alpha_n v_n$. Then
    \[S_{ij}Tv_i = S_{ij}(\alpha_1 v_1 + \cdots + \alpha_n v_n) = \alpha_i v_j,\]
    and
    \[TS_{ij}v_i = T(v_j) = \beta_1 v_1 + \cdots + \beta_n v_n\]
    for some $\beta_1, \ldots, \beta_n \in \f$. Since $S_{ij}T = TS_{ij}$, we have
    \[\alpha_i v_j = \beta_1 v_1 + \cdots + \beta_n v_n.\]
    Since $v_1, \ldots, v_n$ is a basis, it follows that $\beta_k = 0$ for all $k \neq j$ and $\beta_j = \alpha_i$. Therefore, $Tv_j = \alpha_i v_j$. Since $i$ was arbitrary, it follows that $Tv_j = \alpha_j v_j$ for some $\alpha_j \in \f$, hence $T$ sends each basis vector to a scalar multiple of itself.

    Now consider any $\ell \neq j$. Then
    \[S_{\ell j}Tv_\ell = S_{\ell j}(\alpha_\ell v_\ell) = \alpha_\ell v_j,\]
    and
    \[TS_{\ell j}v_\ell = T(v_j) = \alpha_j v_j.\]
    Since $S_{\ell j}T = TS_{\ell j}$, we have $\alpha_\ell v_j = \alpha_j v_j$, which implies $\alpha_\ell = \alpha_j$. Since $\ell$ and $j$ were arbitrary, it follows that all $\alpha_k$ are equal, say $\alpha_k = \lambda$ for all $k$. Therefore, $Tv_k = \lambda v_k$ for all $k$, and hence $T = \lambda I$.
\end{sol}

\begin{exercise}[3a12]
    Suppose $U$ is a subspace of $V$ with $U \neq V$. Suppose $S \in \lin{U, W}$ and $S \neq 0$ (which means that $Su \neq 0$ for some $u \in U$). Define $T: V \to W$ by
    \[Tv = 
    \begin{cases}
        Sv & \text{if } v \in U, \\
        0 & \text{if } v \in V \text{ and } v \notin U. 
    \end{cases}
    \]
    Prove that $T$ is not a linear map on $V$.
\end{exercise}

\begin{sol}
    Since $U \neq V$, there exists $v \in V$ such that $v \notin U$. Let $u \in U$ be such that $Su \neq 0$. Then
    \[T(u + v) = 0,\]
    but
    \[Tu + Tv = Su + 0 = Su \neq 0.\]
    Therefore, $T$ is not linear.
\end{sol}

\begin{exercise}[3a13]
    Suppose $V$ is finite-dimensional. Prove that every linear map on a subspace of $V$ can be extended to a linear map on $V$. In other words, show that if $U$ is a subspace of $V$ and $S \in \lin{U, W}$, then there exists $T \in \lin{V, W}$ such that $Tu = Su$ for all $u \in U$.
\end{exercise}

\begin{sol}
    Let $u_1, \ldots, u_m$ be a basis of $U$. Since $U$ is a subspace of $V$, we can extend this basis to a basis of $V$, say $u_1, \ldots, u_m, v_1, \ldots, v_n$. Define $T: V \to W$ by
    \[Tu = Su \quad \text{for all } u \in U,\]
    and
    \[Tv_k = 0 \quad \text{for all } k = 1, \ldots, n.\]
    Then $T$ is well-defined on the basis of $V$, and we can extend it linearly to all of $V$. For any $u \in U$, we may write
    \[u = \alpha_1 u_1 + \cdots + \alpha_m u_m\]
    for some scalars $\alpha_1, \ldots, \alpha_m \in \f$, and 0 as the coefficients for $v_1, \ldots, v_n$. Then
    \[Tu = \alpha_1 Tu_1 + \cdots + \alpha_m Tu_m + 0 \cdot Tv_1 + \cdots + 0 \cdot Tv_n = \alpha_1 Su_1 + \cdots + \alpha_m Su_m = Su.\]
    Therefore, $T$ extends $S$ to all of $V$.
\end{sol}

\begin{exercise}[3a14]
    Suppose $V$ is finite-dimensional with $\dim V > 0$, and suppose $W$ is infinite-dimensional. Prove that $\lin{V, W}$ is infinite-dimensional.
\end{exercise}

\begin{sol}
    Let $v_1, \ldots, v_n$ be a basis of $V$. Since $W$ is infinite-dimensional, there exists an infinite linearly independent list of vectors $w_1, w_2, \ldots$ in $W$. For each $k \in \n$, define $T_k \in \lin{V, W}$ by
    \[T_k v_j =
    \begin{cases}
        w_k & \text{if } j = 1, \\
        0 & \text{if } j \neq 1.
    \end{cases}\]
    In other words, $T_k$ sends the first basis vector $v_1$ to $w_k$ and all other basis vectors to $0$. We claim that the list $T_1, T_2, \ldots$ is linearly independent in $\lin{V, W}$. Suppose that
    \[a_1 T_1 + a_2 T_2 + \cdots + a_m T_m = 0\]
    for some scalars $a_1, \ldots, a_m \in \f$. Evaluating both sides at $v_1$, we get
    \[a_1 T_1 v_1 + a_2 T_2 v_1 + \cdots + a_m T_m v_1 = a_1 w_1 + a_2 w_2 + \cdots + a_m w_m = 0.\]
    Since $w_1, \ldots, w_m$ is linearly independent in $W$, it follows that $a_1 = \cdots = a_m = 0$. Therefore, $T_1, T_2, \ldots$ is linearly independent in $\lin{V, W}$, and hence $\lin{V, W}$ is infinite-dimensional.
\end{sol}

\begin{exercise}[3a15]
    Suppose $v_1, \ldots, v_m$ is a linearly dependent list of vectors in $V$. Suppose also that $W \neq \set{0}$. Prove that there exist $w_1, \ldots, w_m \in W$ such that no $T \in \lin{V, W}$ satisfies $Tv_k = w_k$ for each $k = 1, \ldots, m$.
\end{exercise}

\begin{sol}
    Since $v_1, \ldots, v_m$ is linearly dependent, there exist scalars $a_1, \ldots, a_m \in \f$, not all zero, such that
    \[a_1 v_1 + \cdots + a_m v_m = 0.\]
    Let $w_1, \ldots, w_m \in W$ be defined by
    \[w_k =
    \begin{cases}
        0 & \text{if } a_k = 0, \\
        w & \text{if } a_k \neq 0,
    \end{cases}\]
    where $w$ is any nonzero vector in $W$. Suppose for contradiction that there exists $T \in \lin{V, W}$ such that $Tv_k = w_k$ for each $k = 1, \ldots, m$. Then
    \[T(a_1 v_1 + \cdots + a_m v_m) = a_1 Tv_1 + \cdots + a_m Tv_m = a_1 w_1 + \cdots + a_m w_m.\]
    Since $a_1 v_1 + \cdots + a_m v_m = 0$, we have $T(0) = 0$, so
    \[0 = a_1 w_1 + \cdots + a_m w_m.\]
    However, by the definition of $w_k$, we have $a_k w_k = a_k w$ for all $k$ such that $a_k \neq 0$, and $a_k w_k = 0$ for all $k$ such that $a_k = 0$. Therefore, the sum $a_1 w_1 + \cdots + a_m w_m$ is a nonzero scalar multiple of $w$, which is nonzero since $w \neq 0$. This is a contradiction. Therefore, no such $T \in \lin{V, W}$ exists.
\end{sol}

\begin{exercise}[3a16]
    Suppose $V$ is finite-dimensional with $\dim V > 1$. Prove that there exist $S, T \in \lin{V}$ such that $ST \neq TS$.
\end{exercise}

\begin{sol}
    Let $v_1, v_2$ be linearly independent vectors in $V$. We may extend this list to a basis of $V$, say $v_1, v_2, \ldots, v_n$. Define $S, T \in \lin{V}$ by
    \[Sv_1 = v_2, \quad Sv_2 = v_1, \quad Sv_k = 0 \text{ for } k > 2,\]
    and
    \[Tv_1 = v_1, \quad Tv_2 = 0, \quad Tv_k = 0 \text{ for } k > 2.\]
    Then
    \[STv_1 = S(Tv_1) = Sv_1 = v_2,\]
    but
    \[TSv_1 = T(Sv_1) = Tv_2 = 0.\]
    Therefore, $ST \neq TS$.
\end{sol}

\begin{exercise}[3a17]
    Suppose $V$ is finite-dimensional. Show that the only two-sided ideals of $\lin{V}$ are $\set{0}$ and $\lin{V}$.
    
    \note{A subspace $\ee$ of $\lin{V}$ is called a two-sided ideal of $\lin{V}$ if $TE \in \ee$ and $ET \in \ee$ for all $E \in \ee$ and all $T \in \lin{V}$.}
\end{exercise}

\begin{sol}
    Let $\ee$ be a two-sided ideal of $\lin V$. If $\ee = \set 0$, then we are done. Otherwise, there exists $E \in \ee$ such that $E \neq 0$. Then there exists $v_1 \in V$ such that $Ev_1 \neq 0$. Since $V$ is finite-dimensional, we may extend the list $\set v_1$ to a basis of $V$, say $v_1, v_2, \ldots, v_n$. Let $w_1 = Ev_1 \neq 0$. We may also extend the list $\set w_1$ to a basis of $V$, say $w_1, w_2, \ldots, w_n$. Define $T_1 \in \lin V$ by
    \[T_1w_1 = v_1, \quad T_1w_k = 0 \text{ for } k > 1.\]
    Then $T_1E \in \ee$ since $\ee$ is a two-sided ideal. Observe that $T_1Ev_1 = T_1(Ev_1) = T_1w_1 = v_1 \neq 0$. Therefore, $T_1E \neq 0$. Consider $S_1 \in \lin V$ defined by
    \[S_1v_1 = v_1, \quad S_1v_k = 0 \text{ for } k > 1.\]
    Then $T_1ES_1 \in \ee$ since $\ee$ is a two-sided ideal. Observe that $T_1ES_1v_1 = T_1Ev_1 = T_1w_1 = v_1 \neq 0$. Therefore, $T_1ES_1 \neq 0$. We now discuss the construction of the identity map $I \in \lin V$.
    
    Construct $T_i, S_i \in \lin V$ for each $i = 1, \ldots, n$ as follows:
    \[T_i w_1 = v_i, \quad T_i w_k = 0 \text{ for } k \neq 1,\]
    and
    \[S_i v_i = v_1, \quad S_i v_k = 0 \text{ for } k \neq i.\]
    Then $T_iES_i \in \ee$ and $T_iES_i v_i = v_i \neq 0$. Since $\ee$ is a subspace, it contains all linear combinations of these maps. In particular, for any $v \in V$, we can write $v = \alpha_1 v_1 + \alpha_2 v_2 + \cdots + \alpha_n v_n$ for some scalars $\alpha_1, \alpha_2, \ldots, \alpha_n \in \f$. Observe that
    \[T_iES_iv = T_iES_i(\alpha_1 v_1) + \cdots + T_iES_i(\alpha_i v_i) + \cdots + T_iES_i(\alpha_n v_n) = \alpha_i v_i.\]
    Consider the map
    \[E^* = T_1ES_1 + T_2ES_2 + \cdots + T_nES_n \in \ee,\]
    where $E^* \in \ee$ since $\ee$ is a subspace. Then
    \[E^*v = T_1ES_1 v + T_2ES_2 v + \cdots + T_nES_n v = \alpha_1 v_1 + \alpha_2 v_2 + \cdots + \alpha_n v_n = v.\]
    Therefore, for any $v \in V$, there exists $E^* \in \ee$ such that $E^*v = v$. Now let $T \in \lin V$ be arbitrary. Then
    \[T = TE^* \in \ee\]
    since $TE^* \in \ee$ and $\ee$ is a two-sided ideal. Therefore, $\ee = \lin V$. We conclude that the only two-sided ideals of $\lin V$ are $\set 0$ and $\lin V$.
\end{sol}